% Evaluation — claude-tunnel bachelor thesis.
\chapter{Evaluation}
\label{ch:evaluation}

Dieses Kapitel beantwortet die Forschungsfragen FF2 (Latenz-Overhead) und FF3
(Netz- und Betriebsartverhalten) empirisch. Es beschreibt das Testbett und die
Methodik (Abschnitt~\ref{sec:eval-methodik}), praesentiert die Ergebnisse
(\ref{sec:eval-ergebnisse}) und ordnet sie samt Limitierungen ein
(\ref{sec:eval-diskussion}). Abschnitt~\ref{sec:eval-security} ergaenzt die
empirische Leistungsbewertung um eine qualitative Sicherheitsbewertung der
Produktivierung aus Kapitel~\ref{ch:produktivierung}.

\section{Testbett und Methodik}
\label{sec:eval-methodik}

Gemessen wird in der Docker-Compose-Topologie aus Kapitel~\ref{ch:architektur}
mit Edge, Agent, Origin und Client auf einem statischen Netz. Link-Stoergroessen
werden per \texttt{tc netem} auf der Edge-Egress konfiguriert (Verzoegerung,
Verlust). Der Client vermisst je Bedingung $n=20$ vollstaendige Roundtrips mit
frisch aufgebauter QUIC-Verbindung; berichtet werden Mittel mit
95\%-\gls{ac:ci}, Standardabweichung sowie p50/p95/p99. Die Matrix umfasst die
Betriebsarten One-shot, Streaming und Datagramm $\times$ Verzoegerung
$\in \{0, 20, 50\}$\,ms $\times$ Verlust $\in \{0, 2\}$\,\%. Aufbau, Messung und
Auswertung sind ausschliesslich containerbasiert und aus dem Repository
reproduzierbar (\texttt{scripts/sweep.sh}, \texttt{plot.sh},
\texttt{tabulate.py}).

\section{Ergebnisse}
\label{sec:eval-ergebnisse}

Tabelle~\ref{tab:latency} fasst alle Bedingungen zusammen.
Abbildung~\ref{fig:vs-delay} zeigt die Skalierung mit der Verzoegerung,
Abbildung~\ref{fig:by-mode} den Vergleich der Betriebsarten und
Abbildung~\ref{fig:vs-loss} den Einfluss des Verlusts.

\begin{table}[t]
  \centering
  \caption{Tunnel-Roundtrip-Latenz unter emulierten Netzbedingungen (\texttt{tc netem}).}
  \label{tab:latency}
  \begin{tabular}{llrrrrrr}
    \toprule
    Verzögerung & Verlust & $n$ & Mittel & p50 & p95 & Min & Max \\
     &  &  & (ms) & (ms) & (ms) & (ms) & (ms) \\
    \midrule
    0ms & 0\% & 5 & 6.8 & 6.6 & 7.9 & 5.8 & 7.9 \\
    0ms & 2\% & 5 & 6.9 & 6.9 & 8.2 & 5.5 & 8.2 \\
    20ms & 0\% & 5 & 92.7 & 88.7 & 109.2 & 88.1 & 109.2 \\
    20ms & 2\% & 5 & 92.2 & 88.2 & 108.5 & 88.0 & 108.5 \\
    50ms & 0\% & 5 & 217.9 & 207.5 & 260.4 & 205.5 & 260.4 \\
    50ms & 2\% & 5 & 237.5 & 207.7 & 306.5 & 207.5 & 306.5 \\
    \bottomrule
  \end{tabular}
\end{table}


\begin{figure}[t]
  \centering
  \includegraphics[width=0.72\textwidth]{latency-vs-delay.png}
  \caption{Tunnel-Roundtrip gegen Edge-Verzoegerung (Modus \texttt{single}), je
  eine Kurve pro Verlust; obere Fehlerbalken bis p95.}
  \label{fig:vs-delay}
\end{figure}

\begin{figure}[t]
  \centering
  \includegraphics[width=0.72\textwidth]{latency-by-mode.png}
  \caption{Vergleich der Betriebsarten (Verlust 0\,\%): single, stream und udp
  ueberlagern sich nahezu.}
  \label{fig:by-mode}
\end{figure}

\begin{figure}[t]
  \centering
  \includegraphics[width=0.72\textwidth]{latency-vs-loss.png}
  \caption{Tunnel-Roundtrip gegen Paketverlust (Modus \texttt{single}), je eine
  Kurve pro Verzoegerung.}
  \label{fig:vs-loss}
\end{figure}

\section{Diskussion und Limitierungen}
\label{sec:eval-diskussion}

\paragraph{FF2 (Latenz-Overhead).} Die Baseline (0\,ms, 0\,\%) liegt ueber alle
Betriebsarten bei rund 8\,ms und misst den Eigenanteil des Tunnels
(QUIC-Handshake, PoW-Rendezvous, Mehrsprung-Relay). Mit der Verzoegerung skaliert
der Roundtrip linear: fuer \texttt{single} gilt naeherungsweise
$\text{RT} \approx 8{,}8 + 6{,}1 \cdot d$ (mit $d$ = einseitige Edge-Verzoegerung
in ms). Der Faktor $\approx 6$ spiegelt wider, dass ein Anwendungs-Roundtrip die
verzoegerte Strecke mehrfach quert. Bei 0\,\% Verlust sind die
\gls{ac:ci} eng ($\pm 2$--$5$\,ms), der Zusammenhang also gut aufgeloest.

\paragraph{FF3 (Netz- und Betriebsartverhalten).} Erstens trifft Paketverlust
fast ausschliesslich den Tail: bei 2\,\% Verlust bleibt der p50 praktisch
unveraendert, waehrend der p99 explodiert (bei 20\,ms von 150 auf 1154\,ms,
Faktor $\approx 7{,}7$). Ursache ist der \gls{ac:pto} von QUIC -- ein waehrend
des Handshakes verlorenes Paket kostet ein RTT-proportionales Timeout. Zweitens
ist die Betriebsart bei 0\,\% Verlust nicht latenzbestimmend
(Abbildung~\ref{fig:by-mode}): single, stream und udp ueberlagern sich, die
Latenz ist verzoegerungsdominiert. Unter Verlust streuen die Punktschaetzer
zwar, doch die 95\%-\gls{ac:ci} ueberlappen deutlich; ein statistisch
belastbarer Modus-Effekt laesst sich aus $n=20$ nicht ableiten.

\paragraph{Limitierungen.} Das Impairment wirkt nur auf der Edge-Egress, nicht
symmetrisch -- die absolute Latenz ist eine untere Schranke. Der Frisch-Dial pro
Iteration ueberzeichnet den Handshake-Anteil bewusst. $n=20$ loest den
Mittelwert bei 0\,\% Verlust eng auf, ist fuer den schweren Tail unter Verlust
aber knapp; die p99-Werte und der Modus-Vergleich unter Verlust sind entsprechend
vorsichtig zu lesen. Die PoW-Schwierigkeit ist in diesem Lauf fixiert; die
Sweep-Achse erlaubt eine gezielte Schwierigkeitsstudie. Echtes NAT-Hole-Punching
bleibt eine Grenze des flachen Container-Testbetts.

\section{Sicherheitsbewertung der Produktivierung}
\label{sec:eval-security}

Die Leistungsmessung erfasst den Eigenanteil des Tunnels, nicht seine
Sicherheitseigenschaften. Da ein quantitatives Angriffs-Benchmark den Rahmen
dieser Arbeit sprengt, werden die in Kapitel~\ref{ch:produktivierung}
eingefuehrten Kontrollen hier qualitativ gegen ein Angreifermodell bewertet.
Tabelle~\ref{tab:security-eval} ordnet jeder Angreiferklasse die wirksame
Kontrolle und das verbleibende Restrisiko zu.

\begin{table}[t]
  \centering
  \caption{Angreifer, wirksame Kontrolle und Restrisiko.}
  \label{tab:security-eval}
  \begin{tabular}{lll}
    \toprule
    Angreifer & Kontrolle & Restrisiko \\
    \midrule
    On-Path-Lauscher & Noise-E2E, TLS~1.3 & Nutzlast keins; Metadaten sichtbar \\
    Rotierter/boeser Edge & interne CA, CA-Root-Trust & keins (Ketten-Validierung) \\
    Unauth. Zugriff & OIDC-RS256 & an Realm-Schluessel gebunden \\
    Rendezvous-Flut & PoW + Rate-Limit & begrenzt, nicht ausgeschlossen \\
    Selbst-Gutschrift & signierter Webhook & keins ohne Provider-Secret \\
    Finanzierter Sybil & -- & ungeloest \\
    \bottomrule
  \end{tabular}
\end{table}

Die zentrale Invariante ist strukturell und nicht politisch: weil der
Noise-Handshake (\cite{perrin2018noise}) Client und Origin direkt
authentifiziert und verschluesselt, sieht der Vermittler nur Chiffrat -- die
Vertraulichkeit der Nutzlast haengt nicht am Wohlverhalten des Betreibers,
sondern an der Kryptografie. Ebenso ersetzt die interne CA das Anheften eines
Blatt-Zertifikats durch Vertrauen in die CA-Wurzel, sodass eine
Zertifikatsrotation die Kette nicht bricht.

Die Verfuegbarkeitskontrollen sind bewusst nur graduell: das
Proof-of-Work-Gatter (\cite{back2002hashcash}) verteuert Fluten, und das
Ratenlimit je Konto verhindert, dass ein einzelnes Konto die Ausgabe
erschoepft; einen \emph{finanzierten} Angreifer, der echte Kosten traegt,
schrecken beide nicht ab. Dies ist die ehrlich benannte offene Flanke. Auch die
Integritaet der Abrechnung ist an ein Geheimnis gebunden: ohne konfiguriertes
Webhook-Secret wird jede Gutschrift abgelehnt (fail-safe), mit falschem Secret
schlaegt die HMAC-Verifikation fehl.

Insgesamt sind Vertraulichkeit, Authentifizierung und Abrechnungsintegritaet
kryptografisch verankert und ohne Zusatzannahmen ueber den Betreiber wirksam,
waehrend Verfuegbarkeit gegen finanzierte Gegner und der Schutz von Metadaten
bewusst ausserhalb des Anspruchs liegen (vgl.\ das Bedrohungsmodell).
